\documentclass[planificacion2024.tex]{subfiles}
\begin{document}
\setstretch{1.5}
\begin{comment}
Elaborar un texto de diagnóstico aulico para alumnos de 6 año de escuela secundaria, el texto debe contener 400 palabras lenguaje apropiado, y comunicar claramente que los alumnos estan en condiciones de realizar el cursado y recalcar siempre que la materia es #Máquinas Eléctricas II. #
    Son alumnos de 6°1° que ya fueron alumnos mios, y que por lo tanto los conozco bien y tambien son excelentes, pero también está la dificultad de que los alumnos no tienen recursos como computadoras y que tampoco podemos acceder a computadoras para poder usar simuladores, ya que no alcanzas las que prestan para usar, y tampoco hay muchos recursos didácticos como fuentes de alimentación y demás qu epuedan servir para la materia. Ademas indicar que la predisposición de los alumnos es muy buena.
\end{comment}

\subsection*{\centering DIAGNÓSTICO AULICO 2024}

\vspace{1cm}

CURSO:6°1°

\vspace{1cm}

Los alumnos de 6°1° año en la materia "Máquinas Eléctricas II" revela un grupo de estudiantes excepcionales que están completamente preparados para cursar y aprovechar al máximo el contenido de la materia. Estos alumnos, pertenecientes al curso 6°1°, ya han sido alumnos míos en el pasado, lo que me permite conocerlos en profundidad y tener plena confianza en sus habilidades y capacidades académicas.

Es importante destacar que estos estudiantes muestran una excelente predisposición para aprender y participar activamente en clase. Su actitud positiva y su interés en la materia son aspectos que contribuyen de manera significativa a un ambiente de aprendizaje estimulante y colaborativo.

Sin embargo, es necesario señalar que enfrentamos ciertas dificultades en relación con los recursos disponibles para la materia. En particular, la falta de acceso a computadoras para accso a simuladores representa un desafío significativo. A pesar de contar con alumnos capacitados y motivados, la ausencia de estos recursos limita nuestras posibilidades de explorar aspectos prácticos y aplicaciones de la teoría en la práctica. Además, la escasez de recursos didácticos, como fuentes de alimentación u otros equipos específicos, dificulta aún más nuestra capacidad para ofrecer una experiencia de aprendizaje completa y enriquecedora.

A pesar de estas limitaciones, confío plenamente en la capacidad de los alumnos para superar los obstáculos y destacarse en la materia. Estoy comprometido a buscar alternativas creativas y soluciones innovadoras para maximizar el potencial de aprendizaje de cada estudiante. Además, cuento con el apoyo y la colaboración de los alumnos y sus familias para superar estos desafíos juntos.

En resumen, los alumnos de 6°1° están listos y preparados para cursar "Máquinas Eléctricas II". Su excelente predisposición, combinada con su habilidad y compromiso, garantiza un semestre exitoso y gratificante. Aunque enfrentamos ciertas dificultades en términos de recursos, estoy seguro de que con trabajo duro, determinación y colaboración, podemos superar cualquier obstáculo y lograr nuestros objetivos educativos.

\end{document}